\documentclass[10pt]{extarticle}
\usepackage[a4paper,margin=1.5cm]{geometry}
\usepackage[T1]{fontenc}
\usepackage{amsmath,amssymb}
\usepackage{graphicx}
\usepackage[export]{adjustbox}
\usepackage{multicol}
\usepackage{enumitem}
\usepackage{booktabs,array,tabularx}
\setlength{\columnsep}{0.7cm}
\setlength{\parindent}{0pt}
\setlength{\parskip}{3pt}
\pagestyle{empty}
\makeatletter
% figure registry (figdefs.tex is generated by the extract step)
\newcommand{\qfigdef}[4]{\@namedef{qf@file@#1}{#2}\@namedef{qf@w@#1}{#3}\@namedef{qf@wide@#1}{#4}}
\newcommand{\qf@check}[1]{\@ifundefined{qf@file@#1}{\PackageError{score2}{Unknown figure id #1}{}}{}}
\newcommand{\qf@img}[4][]{\includegraphics[width=#2\dimexpr\@nameuse{qf@w@#4}pt\relax,max width=#3,max height=0.45\textheight,#1]{\@nameuse{qf@file@#4}}}
% Full-width questions: finalize wraps every question that contains a wide figure in \qfullbegin...\qfullend,
% so the WHOLE question (stem, figures, options) is set across the page in reading order, then the two columns
% resume. The question's head - its number, stem and first figure (\qfullsplit marks the end of the head) -
% is kept on one page, so a number is never left at a page foot with its figure overleaf; the rest may break
% across pages normally, so no big gaps are left. Without a \qfullsplit, the whole question is kept together
% if it is under 45% of a page. Never write these yourself.
\newif\ifqf@full
\newif\ifqf@boxed
\newsavebox{\qf@box}
\newcommand{\qf@open}{\global\qf@fulltrue\par\global\qf@boxedtrue
  \setbox\qf@box\vbox\bgroup\hsize\textwidth\linewidth\textwidth}
\newcommand{\qf@emit}[1]{\par\egroup\global\qf@boxedfalse
  \ifdim\dimexpr\ht\qf@box+\dp\qf@box\relax<#1\textheight\noindent\usebox{\qf@box}\par\else\unvbox\qf@box\fi}
\newcommand{\qfullsplit}{\ifqf@boxed\qf@emit{0.85}\fi}
\newcommand{\qf@close}{\ifqf@boxed\qf@emit{0.45}\fi\global\qf@fullfalse}
\newcommand{\qfullbegin}{\end{multicols}\qf@open}
\newcommand{\qfullend}{\qf@close\begin{multicols}{2}}
\newcommand{\qfullstart}{\qf@open}
\newcommand{\qfulldone}{\qf@close}
% \qfig[scale]{id}: block figure, centred, directly where it is written (size set by the tool).
% Fallback if a wide figure is not inside a full-width question: break out of the columns in place - no floats.
\newcommand{\qfig}[2][1]{\qf@check{#2}%
  \ifqf@full\qf@blk{#1}{#2}\else
  \ifnum\@nameuse{qf@wide@#2}=1 \expandafter\@firstoftwo\else\expandafter\@secondoftwo\fi
  {\end{multicols}\begin{center}\qf@img{#1}{\textwidth}{#2}\end{center}\begin{multicols}{2}}%
  {\qf@blk{#1}{#2}}\fi}
\newcommand{\qf@blk}[2]{\par\smallskip{\centering\qf@img{#1}{\linewidth}{#2}\par}\smallskip}
% \qfigi[scale]{id}: inline figure (inside option lists / tables / minipages); never breaks columns.
\newcommand{\qfigi}[2][1]{\qf@check{#2}\qf@img[valign=c]{#1}{\linewidth}{#2}}
\makeatother
% structure macros - lint reads these to check completeness and figure placement
\newcommand{\papersec}[1]{\par\bigskip{\centering\Large\bfseries PAPER-#1\par}\medskip}
\newcommand{\partsec}[2]{\par\medskip\noindent{\large\bfseries PART-#1: #2}\par\smallskip}
\newcommand{\sect}[1]{\par\smallskip\noindent\textbf{#1}\par}
\newcommand{\qn}[1]{\par\smallskip\noindent\textbf{#1.}~}
\newcommand{\paraq}[2]{\par\smallskip\noindent\textbf{Paragraph for Questions #1 and #2}\par}
\newenvironment{qoptions}{\begin{itemize}[leftmargin=1.8em,itemsep=1pt,topsep=2pt,label={}]}{\end{itemize}}
\newcommand{\opt}[2]{\item[(#1)] #2}
% figure registry (generated by score2_tools.py finalize)
\qfigdef{p1_phy_s0_q01_1}{SRGBT - Test 04 (15 Aug 2021)-img/p1_phy_s0_q01_1.png}{88.1}{0}
\qfigdef{p1_phy_s0_q03_1}{SRGBT - Test 04 (15 Aug 2021)-img/p1_phy_s0_q03_1.png}{109.5}{0}
\qfigdef{p1_phy_s0_q04_1}{SRGBT - Test 04 (15 Aug 2021)-img/p1_phy_s0_q04_1.png}{49.6}{0}
\qfigdef{p1_phy_s0_q05_2}{SRGBT - Test 04 (15 Aug 2021)-img/p1_phy_s0_q05_2.png}{138.3}{0}
\qfigdef{p1_phy_s0_q06_1}{SRGBT - Test 04 (15 Aug 2021)-img/p1_phy_s0_q06_1.png}{89.0}{0}
\qfigdef{p1_phy_s0_q07_1}{SRGBT - Test 04 (15 Aug 2021)-img/p1_phy_s0_q07_1.png}{122.6}{0}
\qfigdef{p1_phy_s0_q08_1}{SRGBT - Test 04 (15 Aug 2021)-img/p1_phy_s0_q08_1.png}{77.1}{0}
\qfigdef{p1_phy_s0_q08_2}{SRGBT - Test 04 (15 Aug 2021)-img/p1_phy_s0_q08_2.png}{48.5}{0}
\qfigdef{p1_phy_s0_q09_1}{SRGBT - Test 04 (15 Aug 2021)-img/p1_phy_s0_q09_1.png}{83.1}{0}
\qfigdef{p1_phy_s0_q10_1}{SRGBT - Test 04 (15 Aug 2021)-img/p1_phy_s0_q10_1.png}{171.8}{0}
\qfigdef{p1_phy_s0_q11_1}{SRGBT - Test 04 (15 Aug 2021)-img/p1_phy_s0_q11_1.png}{89.6}{0}
\qfigdef{p1_phy_s0_q12_1}{SRGBT - Test 04 (15 Aug 2021)-img/p1_phy_s0_q12_1.png}{120.3}{0}
\qfigdef{p1_phy_s0_q13_1}{SRGBT - Test 04 (15 Aug 2021)-img/p1_phy_s0_q13_1.png}{96.3}{0}
\qfigdef{p1_phy_s0_q14_1}{SRGBT - Test 04 (15 Aug 2021)-img/p1_phy_s0_q14_1.png}{134.1}{0}
\qfigdef{p1_phy_s0_q15_1}{SRGBT - Test 04 (15 Aug 2021)-img/p1_phy_s0_q15_1.png}{59.1}{0}
\qfigdef{p1_phy_s0_q16_1}{SRGBT - Test 04 (15 Aug 2021)-img/p1_phy_s0_q16_1.png}{119.5}{0}
\qfigdef{p1_phy_s0_q17_1}{SRGBT - Test 04 (15 Aug 2021)-img/p1_phy_s0_q17_1.png}{99.0}{0}
\qfigdef{p1_phy_s0_q18_1}{SRGBT - Test 04 (15 Aug 2021)-img/p1_phy_s0_q18_1.png}{165.4}{0}
\qfigdef{p1_chem_s0_q22_1}{SRGBT - Test 04 (15 Aug 2021)-img/p1_chem_s0_q22_1.png}{93.2}{0}
\qfigdef{p1_chem_s0_q32_1}{SRGBT - Test 04 (15 Aug 2021)-img/p1_chem_s0_q32_1.png}{175.5}{0}
\qfigdef{p1_chem_s0_q33_1}{SRGBT - Test 04 (15 Aug 2021)-img/p1_chem_s0_q33_1.png}{325.6}{1}
\qfigdef{p1_chem_s0_q35_1}{SRGBT - Test 04 (15 Aug 2021)-img/p1_chem_s0_q35_1.png}{84.5}{0}
\qfigdef{p1_math_s0_q41_1}{SRGBT - Test 04 (15 Aug 2021)-img/p1_math_s0_q41_1.png}{133.9}{0}
\raggedcolumns
\makeatletter\@beginparpenalty=10000\makeatother
\newcolumntype{L}[1]{>{\raggedright\arraybackslash}p{#1}}
\begin{document}
\begin{center}{\Large\bfseries SRGBT --- Test 04 (15 Aug 2021)}\end{center}
\begin{multicols}{2}
\raggedcolumns
\papersec{1}
\partsec{1}{PHYSICS}
\qn{1} A uniform thin hemispherical shell is kept static on an inclined plane of angle $\theta$ = 30$^\circ$ as shown. If the surface of the inclined plane is sufficiently rough to prevent sliding then the angle $\alpha$ made by the plane of hemisphere with inclined plane is :\par
\qfig{p1_phy_s0_q01_1}
\begin{qoptions}
\opt{A}{value of $\mu$ is needed}
\opt{B}{30$^\circ$}
\opt{C}{45$^\circ$}
\opt{D}{60$^\circ$}
\end{qoptions}
\qn{2} A block of mass 2.0 kg is sliding on a frictionless horizontal floor with a velocity of 5.0 m/s due east. It is subjected to a variable net force for 10 sec. resulting in an impulse of 6.0 N.s in the north and 18.0 N.s in the west. What is the total work done by this force ?
\begin{qoptions}
\opt{A}{0.0 J}
\opt{B}{90 J}
\opt{C}{100 J}
\opt{D}{200 J}
\end{qoptions}
\qn{3} In a gravity free space, a hollow cylinder of radius 4R and mass M is being rotated by a constant angular velocity $\omega$ clockwise w.r.t an inertial frame of reference. Another hollow cylinder of radius R and mass m is rotating inside this cylinder, such that the inner surface of larger cylinder without slipping and its centre of mass appears stationary from inertial frame Dynamic friction coefficient between cylinders is $\mu$. If the normal reaction between cylinders is denoted by N.\par
\qfig{p1_phy_s0_q03_1}
A small tangential velocity v is given to the smaller cylinder towards right. What will be N, once the steady state is achieved. (Outer cylinder is rotating about fixed axis passing through center O and perpendicular to the plane of paper).
\begin{qoptions}
\opt{A}{m$\dfrac{\mathrm v^2}{16\mathrm R}$}
\opt{B}{m$\dfrac{\mathrm v^2}{12\mathrm R}$}
\opt{C}{m$\dfrac{(3\mathrm R\cdot(3\omega/4)+\mathrm v)^2}{3\mathrm R}$}
\opt{D}{zero}
\end{qoptions}
\qn{4} The given system is released from rest then find change in acceleration of particle `P' in the interval it just comes in contact with the pulley to just before it leaves the pulley. Pulley and string are massless but sufficiently rough :\par
\qfig{p1_phy_s0_q04_1}
\begin{qoptions}
\opt{A}{$\dfrac{2\mathrm g}3\sqrt{1+(2+\pi)^2}$}
\opt{B}{$\dfrac{\mathrm g}3\sqrt{1+(2+\pi)^2}$}
\opt{C}{$\dfrac{\mathrm g}3\sqrt{1+(1+\pi)^2}$}
\opt{D}{$\dfrac{\mathrm g}3\sqrt{1+(1+\pi)^2}$}
\end{qoptions}
\qn{5} Two identical particles connected with a massless rod are placed on a frictionless horizontal floor, which terminates into a uniform frictionless slope of inclination $\theta$ = 45$^\circ$. Find minimum velocity given to the particles towards the slope so that the rod can rise completely on the slope. Collisions of the particle with the slope are perfectly inelastic.\par
\qfig{p1_phy_s0_q05_2}
\begin{qoptions}
\opt{A}{$\sqrt{\dfrac{\mathrm g\ell}{2\sqrt2}}$}
\opt{B}{$\sqrt{\dfrac{\sqrt2\mathrm g\ell}3}$}
\opt{C}{$\sqrt{\dfrac{2\sqrt2\mathrm g\ell}3}$}
\opt{D}{$\sqrt{\dfrac{3\mathrm g\ell}{2\sqrt2}}$}
\end{qoptions}
\qn{6} Two bodies A and B of masses 5.00 kg and 10.0 kg moving in free space in opposite directions with velocities 4.00 m/s and 0.50 m/s respectively undergo a head-on collision. The force F of their mutual interaction varies with time t according to the given graph. What can you conclude from the given information ?\par
\qfig{p1_phy_s0_q06_1}
\begin{qoptions}
\opt{A}{Period of deformation is 0.20 s.}
\opt{B}{Coefficient of restitution is 0.5.}
\opt{C}{Body A will move with velocity 0.5 m/s in the original direction.}
\opt{D}{Body B will move with velocity 1.75}
\end{qoptions}
\qn{7} A box of mass 4 kg placed on a horizontal floor, is acted upon by a horizontal force F that varies with time t as shown in the given graph. If coefficients of static and kinetic frictions are 0.25, on which of the following conclusions can you arrive ?\par
\qfig{p1_phy_s0_q07_1}
\begin{qoptions}
\opt{A}{The box starts moving at the instant t = 2.0 s.}
\opt{B}{The maximum velocity acquired by the box is 5.0 m/s}
\opt{C}{The box stops at the instant t = 9.0 s.}
\opt{D}{Modulus of average power of the force F is more than that of the frictional force.}
\end{qoptions}
\qn{8} A uniform solid cylinder of mass m, radius R is at rest on an extremely rough horizontal surface. Now a horizontal force F = kt where k = constant and t = time, is applied at the highest point on the cylinder. Assume that the cylinder is not sliding at all.\par
\qfig{p1_phy_s0_q08_1}
\begin{qoptions}
\opt{A}{The friction force acting on the cylinder varies with time as \qfigi{p1_phy_s0_q08_2}}
\opt{B}{Velocity of the highest point after time t will be $\dfrac{2\mathrm{kt}^2}{3\mathrm m}$}
\opt{C}{Speed of the particle which is at highest point of the cylinder after time t will be $\dfrac{4\mathrm{kt}^2}{3\mathrm m}$}
\opt{D}{If the coefficient of friction between the ground and the cylinder is $\mu$, the cylinder will start sliding at t = $\dfrac{3\mu\mathrm{mg}}{\mathrm k}$}
\end{qoptions}
\qn{9} 4 identical balls of radius R and mass m are lying in a gravity free space. The balls are in contact and there centres are forming vertices of a square of side 2R in horizontal plane. One identical ball travelling vertically with a speed v hits the four balls symmetrically. The collision is perfectly elastic. The centre of the 4 stationary ball are at (R, R, 0), ($-$R, R, 0), (R, $-$R, 0) and ($-$R, $-$R, 0) initially.\par
\qfig{p1_phy_s0_q09_1}
\begin{qoptions}
\opt{A}{The speed of ball 1 after collision is $\dfrac{\mathrm v}{\sqrt2}$}
\opt{B}{The speed of ball 1 after collision is $\dfrac{\sqrt2}3$v}
\opt{C}{The speed of ball which was initially travelling with speed v just after collision is v/3}
\opt{D}{The speed of ball which was initially travelling with speed v just after collision is v/$\sqrt3$}
\end{qoptions}
\qn{10} A 60 mm radius wheel is connected to a fixed support D by two links AB and BD. At the instant shown, the velocity of the centre A of the wheel is 30 m/s to the left. Rod AB and BD have equal length $\ell$ = 5m. Then at the given instant :\par
\qfig{p1_phy_s0_q10_1}
\begin{qoptions}
\opt{A}{The angular velocity of rod AB will be 3$\sqrt2$ rad / sec}
\opt{B}{The angular velocity of rod BD is 3$\sqrt2$ rad / sec}
\opt{C}{Velocity of pin B will be 5$\sqrt2$ m / s}
\opt{D}{Velocity of pin B will be 15$\sqrt2$ m / s}
\end{qoptions}
\qn{11} The given uniform disc is rotating about vertical y-axis with constant angular speed $\omega_0$. If surface is sufficiently rough then (OA is a massless rod and plane of disc is perpendicular to the rod\par
\qfig{p1_phy_s0_q11_1}
\begin{qoptions}
\opt{A}{angular velocity of disc about an axis passing through the centre O and perpendicular to the disc is 2$\omega_0$.}
\opt{B}{Kinetic energy of the disc is $\dfrac{25}8\mathrm{mR}^2\omega_0^2$}
\opt{C}{Acceleration of point of contact with ground is 2$\sqrt5\omega_0^2$R}
\opt{D}{Kinetic energy in the reference frame of center of mass `O' is $\dfrac98\mathrm{mR}^2\omega_0^2$}
\end{qoptions}
\qn{12} The given uniform rod is hinged at point A and rotates with constant angular velocity $\omega$ about vertical axes AB.\par
\qfig{p1_phy_s0_q12_1}
\begin{qoptions}
\opt{A}{Angular momentum of the rod about point A is $\dfrac{\mathrm m\ell^2\omega}{12}$}
\opt{B}{Angular momentum of the rod about point A is $\dfrac{\mathrm m\ell^2\omega}6$}
\opt{C}{Angular momentum of the rod about axis AB is $\dfrac{\mathrm m\ell^2\omega}{12}$}
\opt{D}{Angular momentum of the rod and torque about point A are always perpendicular to each other.}
\end{qoptions}
\qn{13} The given hollow sphere of mass M and radius R is in pure rolling with constant linear velocity v = $\sqrt{\mathrm{gR}}$ . Then at the given instant :\par
\qfig{p1_phy_s0_q13_1}
\begin{qoptions}
\opt{A}{Force exerted by the upper hemisphere on the lower hemisphere is zero}
\opt{B}{Force exerted by the upper hemisphere on the lower hemisphere is Mg/4}
\opt{C}{Kinetic energy of lower hemisphere is $\dfrac13$Mv$^2$}
\opt{D}{Kinetic energy of lower hemisphere is $\dfrac16$Mv$^2$}
\end{qoptions}
\qn{14} Three identical uniform cylinders are mounted on frictionless horizontal axles passing through their centres. They are given equal angular velocities $\omega_0$ each and then gradually brought into contact preventing any vertical movement and rotation of their axles. The angular velocities of the cylinders after they stop slipping is $\omega_0$x then find x ?\par
\qfig{p1_phy_s0_q14_1}
\qn{15} Inside a fixed frictionless spherical cavity of radius r, a dumbbell is held touching the bottom of the cavity at its lower end. The dumbbell consists of two identical small balls each of mass m affixed at the ends of a light rod of length r$\sqrt2$ . The force of interaction between the lower ball and the cavity be immediately after the dumbbell is released is x mg then find x.\par
\qfig{p1_phy_s0_q15_1}
\qn{16} The wheel of radius R rolls without slipping on horizontal rough surface, and its centre O has an horizontal acceleration a$_0$ in forward direction. A point P on the wheel is a distance r from O and angular position $\theta$ from horizontal. For a given values of a$_0$, R and r, the angle $\theta$ for which point P has no acceleration in this position is $\sin^{-1}\left(\dfrac{\mathrm{xr}}{\mathrm{yR}}\right)$ then find x + y ? (x and y are lowest possible integers).\par
\qfig{p1_phy_s0_q16_1}
\qn{17} The given disc is slightly rotated and then released. The angular acceleration of the system after the particle `A' has rotated by angle 90$^\circ$ about centre O is $\dfrac{\mathrm{xg}}{\mathrm{yR}}$ then find x + y ? The surface is sufficiently rough. (x and y are lowest possible integers).\par
\qfig{p1_phy_s0_q17_1}
\qn{18} A block of mass m can slide without friction on a horizontal floor. a small ball of same mass m is suspended by a light inextensible cord of length $\ell$ from a hanger fixed at the top of a vertical pole mounted on the block so that the ball can swing freely in a vertical plane. The ball is released from a horizontal position as shown. The maximum tensile force in the string during the subsequent motion is x mg then find x ?\par
\qfig{p1_phy_s0_q18_1}
\raggedcolumns
\papersec{1}
\partsec{2}{CHEMISTRY}
\qn{19} Translational, rotational and vibrational degrees of freedom contribute to the molar heat capacity at constant volume, C$_\mathrm V$. For AX$_{3(\mathrm g)}$, C$_\mathrm V$ is 18 cal/deg and for M$_2$H$_{2(\mathrm g)}$ C$_\mathrm V$ is 19 cal deg$^{-1}$, one can therefore conclude that :
\begin{qoptions}
\opt{A}{Both AX$_3$ and M$_2$H$_2$ are linear}
\opt{B}{Both AX$_3$ and M$_2$H$_2$ are non-linear}
\opt{C}{AX$_3$ is linear and M$_2$H$_2$ is non-linear}
\opt{D}{AX$_3$ is non-linear and M$_2$H$_2$ is linear}
\end{qoptions}
\qn{20} Which of the following option is incorrect ?
\begin{qoptions}
\opt{A}{The difference between ionisation enthalpy ($\Delta$H$_\mathrm{IE}$) and ionisation energy (IE) is directly depencent on absolute temperature.}
\opt{B}{heat added to the system at lower temperature causes greater randomness as compared to same amount of heat added at higher temperature.}
\opt{C}{A container having N$_2$O$_4$ gas if heated will show an increase in intensity of brown colour.}
\opt{D}{It is easier to remove H$^+$ from HPO$_4^{2-}$ as compared to that from H$_2$PO$_4$.}
\end{qoptions}
\qn{21} From the following data.\par
C(s, graphite) + O$_2$(g) $\longrightarrow$ CO$_2$(g) \hfill $\Delta$H$^\circ_\mathrm{reaction}$ = $-$393.5 kJ / mole\par
H$_2$(g) + $\dfrac12$O$_2$(g) $\longrightarrow$ H$_2$O($\ell$) \hfill $\Delta$H$^\circ_\mathrm{reaction}$ = $-$286 kJ / mole\par
2C$_2$H$_6$(g) + 7O$_2$(g) $\longrightarrow$ 4CO$_2$(g) + 6H$_2$O($\ell$) \hfill $\Delta$H$^\circ_\mathrm{reaction}$ = $-$3120 kJ / mole\par
Calculate the standard enthalpy of formation of C$_2$H$_6$ (g) (in kJ/mole).
\begin{qoptions}
\opt{A}{$-$85}
\opt{B}{$-$170}
\opt{C}{170}
\opt{D}{85}
\end{qoptions}
\qn{22} If four identical samples of an ideal gas initially at similar state (P$_0$, V$_0$, T$_0$) are allowed to expand to double their volumes byy four different process :\par
(I) by isothermal irrevrsible process.\par
(II) by reversible process such that P$^2$V = constant\par
(III) by reversible adiabatic process.\par
(IV) by irreversible adiabatic expansion against constnt external pressure.\par
Then in the graph shown the final state is represented by four different points then, the correct match can be :\par
\qfig{p1_chem_s0_q22_1}
\begin{qoptions}
\opt{A}{1--(I), 2--(II), 3--(III), 4--(IV)}
\opt{B}{1--(II), 2--(I), 3--(IV), 4--(III)}
\opt{C}{2--(III), 3--(II), 4--(I), 1--(IV)}
\opt{D}{3--(II), 1--(I), 2--(IV), 4--(III)}
\end{qoptions}
\qn{23} Which of the following options represents correct sequence of True(T) or False(F) statements.\par
\textbf{Statement - I} No Gibbs free energy change occurs four vapourisation at triple point.\par
\textbf{Statement - II} Entropy driven processes are spontaneous at high temperature.\par
\textbf{Statement - III} $\Delta$G$^\circ$ of a reaction will increase with increase in temperature if the reaction is endothermic.
\begin{qoptions}
\opt{A}{T T T}
\opt{B}{T T F}
\opt{C}{F T F}
\opt{D}{T F F}
\end{qoptions}
\qn{24} Mention the correct statement(s) among the following :
\begin{qoptions}
\opt{A}{Keeping initial states same, final temperature for reversible adiabatic expansion is more than irreversible adiabatic expansion if final volume is also same.}
\opt{B}{A(s) $\rightleftharpoons$ A(g) \quad $\log_{10}=4-\dfrac{2000}{\mathrm T}$ \quad $\Delta$H$_\mathrm{Sub}$ for A (s) is 9.212 kCal/mole.}
\opt{C}{For 1 mole of N$_2$ gas State A $\xrightarrow{\text{Single step Irreversible}}$ State B (State A : 0$^\circ$C, 1 atm; State B : 0$^\circ$C, 2 atm) Q$_\mathrm{Process}$ = $-$22.4 L-atm if P$_\mathrm{ext}$ = 2 atm}
\opt{D}{For 1 mole of O$_2$ gas State A $\xrightarrow{\text{Single step Irreversible}}$ State B (State A : 0$^\circ$C, 1 atm; State B : 0$^\circ$C, 2 atm) $\Delta$H = 0 and W$_\mathrm{process}$ = 44.8 L-atm if P$_\mathrm{ext}$ = 2 atm}
\end{qoptions}
\qn{25} One mole of an ideal gas (C$_\mathrm{V,m}=\dfrac52$R ) at 300 K and 5 atm is expanded adiabatically against constant atmospheric pressure to a final pressure of 2 atm. Then which is/are true.
\begin{qoptions}
\opt{A}{Final temperature of the gas is approx 251.4 K.}
\opt{B}{Final temperature of the gas is approx 270 K.}
\opt{C}{P$_1$V$_1^\gamma$ = P$_2$V$_2^\gamma$ is valid for the process.}
\opt{D}{If the process is carried out in infinite steps then final temperature will be less than the temperature observed in above process.}
\end{qoptions}
\qn{26} If an equation like N$_2$O$_4$(g) $\rightleftharpoons$ 2NO$_2$(g) , the $\Delta$G (reaction) expression given is\par
$\Delta$G = 4.73 kJ mol$^{-1}$ + 2.5 kJ mol$^{-1}$ $\ell$n$\left\{\dfrac{\mathrm P_{\mathrm{NO}_2}^2}{\mathrm P_{\mathrm N_2\mathrm O_4}}\right\}$ then
\begin{qoptions}
\opt{A}{If small amount of N$_2$O$_4$ is introduced in a closed container, then reaction moves in forward direction spontaneously.}
\opt{B}{Initially, $\Delta$G$^\circ$ value is negative for the given decomposition.}
\opt{C}{At equilibrium, Q$_\mathrm P$ = K$_\mathrm P$ = antilog $\left|\dfrac{-\Delta\mathrm G^\circ}{\mathrm{RT}}\right|$}
\opt{D}{At equilibrium state $\Delta$G must be zero.}
\end{qoptions}
\qn{27} Select the correct option(s) :
\begin{qoptions}
\opt{A}{$\Delta$H$^\circ_\mathrm f$[H$_2$O($\ell$)] is zero at 25$^\circ$C}
\opt{B}{$\Delta$H$^\circ_\mathrm f$[O$_2$(g)] is zero at 25$^\circ$}
\opt{C}{S$^\circ$[O$_2$(g)] is zero at 25$^\circ$C}
\opt{D}{Entropies of aqueous ions can be negative or positive.}
\end{qoptions}
\qn{28} Which of the following statement are true?
\begin{qoptions}
\opt{A}{The entropy of the universe increases and tends towards the maximum value}
\opt{B}{All natural processes are irreversible}
\opt{C}{Entropy of the system is maximum when the system is at equilibrium}
\opt{D}{According to the third law of thermodynamics, S$^0_\mathrm{0K}$ = 0 for any element or compound.}
\end{qoptions}
\qn{29} 1 mole of a real gas changes its state from -- A (2 bar, 3L, 100 K) to state -- B (2 bar, 5 L, 200 K) at constant pressure and finally to state -- C (3 bar, 10 L, 300 K). If $\Delta$U$_\mathrm{BC}$ = 110 J and C$_\mathrm{p,m}$of gas = 3 R = 3 $\times$ 8.3 JK$^{-1}$ mol$^{-1}$ then choose the correct option(s) :
\begin{qoptions}
\opt{A}{W$_\mathrm{AB}$ = 830 J}
\opt{B}{$\Delta$H$_\mathrm{AC}$ = 4600 J}
\opt{C}{$\Delta$U$_\mathrm{AC}$ = 2200 J}
\opt{D}{$\Delta$U$_\mathrm{AC}$ = 1770 J}
\end{qoptions}
\qn{30} Assume ideal gas behaviour for all the gases considered \& vibrational degree of freedom to be active separate equimolar samples of He, H$_2$, SO$_2$ \& CH$_4$ were subjected to a two step process as mentioned. Initially all are at same state of temperature \& pressure.\par
\textbf{Step I} $\to$ All undergo reversible adiabatic expansion to attain same final volume, which is double the original volume thereby causing the decreases in their temperature.\par
\textbf{Step II} $\to$ After step I all are given appropriate amount of heat isochorically to restore the original temperature.\par
Mark the correct option(s).
\begin{qoptions}
\opt{A}{Due to step I only, the decreases in temperature will be maximum for CH$_4$.}
\opt{B}{During step II, heat given will be maximum for CH$_4$.}
\opt{C}{There will be no change in internal energy for any of the gas after both the steps of process are completed.}
\opt{D}{The P--V graph of CH$_4$ \& SO$_2$ will be same.}
\end{qoptions}
\qn{31} Which is correct relation between various enthalpies given all enthalpies are measured under similar pressure and temperature conditions-
\begin{qoptions}
\opt{A}{2$\Delta_\mathrm f$H$^\circ$(HCl,g) = BE$_\mathrm{(H-H)}$ + BE$_\mathrm{(Cl-Cl)}$ $-$ 2BE$_\mathrm{(H-Cl)}$}
\opt{B}{$\Delta_\mathrm f$H$^\circ$(CO$_2$,g) = $\Delta$H$^\circ_\mathrm{comb.}$ (C, graphite)}
\opt{C}{$\Delta$H$^\circ$ for conversion of d-lactic acid to $l$-lactic acid is zero}
\opt{D}{$\Delta$H$^\circ_\mathrm{neutralization}$ (HCl + NaOH) is equal to $\Delta$H$^\circ$ for the reaction H$^+$(aq.) + OH$^-$(aq.) $\to$ H$_2$O($l$)}
\end{qoptions}
\qn{32} One mole of a monoatomic ideal gas is subjected to a process as shown. Calculate the maximum temperature obtained during the process (in Kelvin)\par
\qfig{p1_chem_s0_q32_1}
\textbf{[If T$_\mathrm{max}$ = $\dfrac{\mathrm X}{\mathrm R}$, then fill X in the OMR sheet]}
\qfullbegin
\qn{33} Calculate $\Delta$H$_\mathrm{resonance\;energy}$ of pyridine.\par
Given data : $\Delta$H$^\circ_\mathrm f$ pyridine$_{(\ell)}$ = $-$175 kJ/mole\par
\qfig{p1_chem_s0_q33_1}
\qfullend
\qn{34} Certain quantity of an ideal gas is expanded from 20 L to 60 L isothermally in three steps. Pressure at these three stages are 1.5 atm, 1.2 atm and 1 atm respectively now the gas is brought back to the initial condition by these three stages again. Calculate the total work done on the system during this cycle in joule.\par
(Take : 1 atm $-$ L = 100 J)
\qn{35} An ideal gas is carried through a thermodynamic cycle ABCD.\par
\qfig{p1_chem_s0_q35_1}
P$_2$ = 10 bar \hfill V$_1=\dfrac1{\ell\mathrm n10}$ litre\par
P$_1$ = 1 bar \hfill V$_2=\dfrac4{\ell\mathrm n10}$ litre\par
Consisting of two isobaric and two isothermal process as shown in figure. Calculate the magnitude of net work in the entire cycle in litre bar.
\qn{36} A reversible cyclic process involves 6 steps. In step -1, 3 system absorb 500 J, 800 J of heat from a heat reservoir at temperature 250K \& 200K respectively. Step 2, 4, 6, are adiabatic such that the temperature of one reservoir changes to that of next. Total work done by the system in whole cycle is 700 J. Find the temperature during step 5 if it exchange heat from a reservoir at temperature T$_5$
\raggedcolumns
\papersec{1}
\partsec{3}{MATHEMATICS}
\qn{37} Circles are drawn through the points (a, b) and (b, $-$a) such that the chord joining the two points subtends an angle of 45$^\circ$ at any point of the circumference. The the distance between the centres is :
\begin{qoptions}
\opt{A}{$\sqrt3$ times the radius of either circle}
\opt{B}{2 times the radius of either circle}
\opt{C}{$\dfrac1{\sqrt2}$ times the radius of either circle}
\opt{D}{$\sqrt2$ times the radius of either circle}
\end{qoptions}
\qn{38} If four points P, Q, R, S in the plane be taken and the square of the length of the tangent from P to the circle on QR as diameter be denoted \{P, QR\}, then :
\begin{qoptions}
\opt{A}{\{P, RS\} $-$ \{P, QS\} + \{Q, PR\} $-$ \{Q, RS\} = 0}
\opt{B}{\{P, RS\} + \{P, QS\} + \{Q, PR\} + \{Q, RS\} = 0}
\opt{C}{\{P, RS\} $-$ \{P, QS\} + \{Q, PR\} $-$ \{Q, RS\} = 1}
\opt{D}{\{P, RS\} $-$ \{Q, RS\} $-$ \{Q, PR\} $-$ \{P, QS\} = 0}
\end{qoptions}
\qn{39} A variable triangle ABC is circumscribed about a fixed circle of unit radius, Side BC always touches the circle at fixed point D. If B and C vary in such a way that BD.DC = 2, then the locus of vertex A will lie on straight line:
\begin{qoptions}
\opt{A}{parallel to side BC}
\opt{B}{perpendicular to side BC}
\opt{C}{Making an angle of $\dfrac\pi6$ with side BC}
\opt{D}{Making an angle of $\sin^{-1}\left(\dfrac23\right)$ with side BC}
\end{qoptions}
\qn{40} Let A and B be two fixed points and P, another point in the plane, moves in such a way that k$_1$PA + k$_2$PB=k$_3$, where k$_1$, k$_2$ and k$_3$ are real constants. Then which of the following is locus of P :
\begin{qoptions}
\opt{A}{an ellipse, if k$_2$ = $-$1 and k$_1$, k$_3$ $>$ 0}
\opt{B}{a circle, if k$_1$ $>$ 0 , k$_2$ $<$ 0 , k$_3$ = 0}
\opt{C}{an ellipse, if k$_1$ = k$_2$ $>$ 0 and k$_3$ $>$ 0}
\opt{D}{a circle, if k$_1$ = 0 and k$_2$, k$_3$ $>$ 0}
\end{qoptions}
\qn{41} adjacent figure represents a equilateral triangle ABC of side length 2 units. Locus of vertex `C' ass the side AB slides along the coordinates axes is :\par
\qfig{p1_math_s0_q41_1}
\begin{qoptions}
\opt{A}{x$^2$ + y$^2$ $-$ xy + 1 = 0}
\opt{B}{x$^2$ + y$^2$ + xy$\sqrt3$ = 1}
\opt{C}{x$^2$ + y$^2$ = 1 + xy$\sqrt3$}
\opt{D}{x$^2$ + y$^2$ $-$ xy$\sqrt3$ + 1 = 0}
\end{qoptions}
\qn{42} Which of the following statementss is true about circles x$^2$ + y$^2$ + 2x = 0 and x$^2$ + y$^2$ $-$ 6x = 0 :
\begin{qoptions}
\opt{A}{length of the transverse common tangents is $\sqrt2$ units}
\opt{B}{length of the direct common tangents is 2$\sqrt3$ units}
\opt{C}{length of the transverse common tangents is 0 units}
\opt{D}{area of trapezium, enclosed by the common tangents and chords of contact is 6$\sqrt3$ units$^2$}
\end{qoptions}
\qn{43} A line segment of length 7 units has one of its ends on x-axis and other end on y-axis. A point P divides this line segment in the ratio 3 : 4 locus of P will be an ellipse with :
\begin{qoptions}
\opt{A}{length of latus rectum $\dfrac92$}
\opt{B}{eccentricity $\dfrac{\sqrt7}4$}
\opt{C}{length of major axis 8}
\opt{D}{None of these}
\end{qoptions}
\qn{44} A circle is touching the lines x $-$ 2y = 3 and 2x + y = 1, the locus of its centre is :
\begin{qoptions}
\opt{A}{3x $-$ y $-$ 4 = 0}
\opt{B}{x $-$3y $-$ 4 = 0}
\opt{C}{x + 3y + 2 = 0}
\opt{D}{3x + y $-$ 2 = 0}
\end{qoptions}
\qn{45} If the line x + 2y = 3 is rotated about its point of intersection with the line x = y by an angle of 45$^\circ$, then it's equations in the new position is/are :
\begin{qoptions}
\opt{A}{3x $-$ y $-$ 2 = 0}
\opt{B}{x + 3y $-$ 4 = 0}
\opt{C}{3x + y $-$ 4 = 0}
\opt{D}{x $-$ 3y + 2 = 0}
\end{qoptions}
\qn{46} A circle S$_2$ = 0 has its centre lying on the circle S$_1$ : x$^2$ + y$^2$ = 4 and radius half of that of S$_1$. If a common tangent of these circles is inclined at an angle of 30$^\circ$ with positive direction of x-axis, then possible equation(s) of S$_2$ is/are?
\begin{qoptions}
\opt{A}{x$^2$ + y$^2$ + 4x + 3 = 0}
\opt{B}{x$^2$ + y$^2$ $-$ 4x + 3 = 0}
\opt{C}{x$^2$ + y$^2$ $-$ 2x $-$ 2$\sqrt3$ y + 3 = 0}
\opt{D}{x$^2$ + y$^2$ + 2x + 2$\sqrt3$ y + 3 = 0}
\end{qoptions}
\qn{47} Cosider the points O(0, 0), A(0, 1), B(1, 1) in the xy plane, suppose that points C(x, 1) and D(1, y) on a line passing through O are chosen such that 0 $<$ x $<$ 1. Let sum of areas of DAC and DBCD be `S' then which of the following are correct?
\begin{qoptions}
\opt{A}{min. value of S is irrational in $\left(\dfrac13,\dfrac12\right)$}
\opt{B}{min. value of S is irrational in $\left(\dfrac23,1\right)$}
\opt{C}{The value of x for minimum value of S lies in $\left(\dfrac23,1\right)$}
\opt{D}{The value of x for minimum values of S lies in $\left(\dfrac13,\dfrac12\right)$}
\end{qoptions}
\qn{48} A tangent drawn from the point (4, 0) to circle x$^2$ + y$^2$ = 8 touches it at a point A in the first quadrant. The coordinates of another point B on the circle such that $\ell$(AB) = 4 :
\begin{qoptions}
\opt{A}{(2, $-$2)}
\opt{B}{($-$2, 2)}
\opt{C}{($-$2, $-$2)}
\opt{D}{(2, 2)}
\end{qoptions}
\qn{49} Two circles of equal radius, having centers at (2, 3) and (5, 6) cut each other orthogonally. Their common chord will pass through the point(s)
\begin{qoptions}
\opt{A}{(2, 6)}
\opt{B}{(6, 4)}
\opt{C}{(7, 1)}
\opt{D}{(3, 5)}
\end{qoptions}
\qn{50} The vertices of a triangle is (4, 3), ($-$4, 3) and (5, 0). The orthocenter of the triangle is (a, b). The $|$a $-$b$|$ is equal to :
\qn{51} A ray along line x + 2y $-$ 3 = 0 is reflected by mirror x = y, if reflected line has equation ax + by $-$ 3 = 0. Find a+b.
\qn{52} In the triangle ABC coordinates of A are (1, 2), Bisector of internal angle B is 2x $-$ y = 5 and perpendicular bisector of the side AC is 3x $-$ 4y = 20. If equation of BC is ax + by = c, then $\dfrac{\mathrm c-\mathrm a}{\mathrm b}$=?
\qn{53} The chord of contact of tangents from A to the circle x$^2$ + y$^2$ =a$^2$ passes through a point B. If 3 and 4 are the length of tangents from A and B on the given circle then AB equals to :
\qn{54} P is a point inside the triangle ABC. Lines are drawn through P, parallel to the sides of the triangle. The three resulting triangles with the vertex at P have areas 4, 9 and 49 sq. units. The area of the triangle ABC is
\end{multicols}
\end{document}
